\documentclass[aspectratio=169]{beamer}

% Theme selection
\usetheme{Madrid}
\usecolortheme{beaver}

% Packages
\usepackage{graphicx}
\usepackage{amsmath}
\usepackage{algorithm}
\usepackage{algpseudocode}
\usepackage{listings}
\usepackage{booktabs}
\usepackage{hyperref}

% Code snippet styles
\lstset{
    language=C++,
    basicstyle=\ttfamily\footnotesize,
    keywordstyle=\color{blue}\bfseries,
    commentstyle=\color{green!50!black},
    stringstyle=\color{red},
    numbers=left,
    numberstyle=\tiny\color{gray},
    stepnumber=1,
    showstringspaces=false,
    tabsize=2,
    frame=single
}

% Title Slide Info
\title[Parallel N-Body Simulation]{Astrophysics N-Body Simulation}
\subtitle{Parallelizing Gravitational Interactions via Direct and Barnes-Hut Methods}
\author{Your Name}
\institute{Your University / Organization}
\date{\today}

\begin{document}

% Slide 1: Title
\begin{frame}
    \titlepage
\end{frame}

% Slide 2: Outline
\begin{frame}{Outline}
    \tableofcontents
\end{frame}

\section{Introduction \& Motivation}
% Slide 3
\begin{frame}{What is the N-Body Problem?}
    \begin{itemize}
        \item Predicting the individual motions of a group of celestial objects interacting with each other gravitationally.
        \item Inspired by Newton's law of universal gravitation.
        \item Historically applied to the Sun, Earth, and Moon.
        \item Currently applied to entire star clusters, galaxies, and cosmological structure formation.
    \end{itemize}
\end{frame}

% Slide 4
\begin{frame}{Applications of N-Body Simulations}
    \begin{itemize}
        \item \textbf{Astrophysics:} Galaxy collisions, dark matter halos, star cluster dynamics.
        \item \textbf{Molecular Dynamics:} VdW forces between atoms behave similarly to gravity (inverse square laws).
        \item \textbf{Plasma Physics:} Charged particle interactions.
        \item \textbf{Computer Graphics:} Particle systems for visual effects (smoke, explosions).
    \end{itemize}
\end{frame}

% Slide 5
\begin{frame}{The Computational Challenge}
    \begin{itemize}
        \item Gravity has infinite range: \textit{every} particle interacts with \textit{every other} particle.
        \item For $N$ bodies, there are $\frac{N(N-1)}{2}$ unique interactions.
        \item The computational complexity is $\mathcal{O}(N^2)$.
        \item A simulation with 1 billion stars requires $10^{18}$ calculations per time step.
        \item \textbf{Conclusion:} Brute-force is not scalable for real astrophysics.
    \end{itemize}
\end{frame}

% Slide 6
\begin{frame}{Project Objectives}
    \begin{enumerate}
        \item Build a custom simulation engine in C++ from scratch.
        \item Implement the \textbf{Direct $\mathcal{O}(N^2)$} computation method.
        \item Implement the optimized \textbf{Barnes-Hut $\mathcal{O}(N \log N)$} tree algorithm.
        \item Utilize \textbf{OpenMP} for shared-memory multi-core parallelization.
        \item Apply a \textbf{Leapfrog Integrator} for perfect energy conservation.
        \item Provide Python data visualizations of physical scaling and orbital tracks.
    \end{enumerate}
\end{frame}

\section{Theoretical Framework}
% Slide 7
\begin{frame}{Governing Equations}
    Newton's law states the force $\mathbf{F}_{ij}$ exerted on body $i$ by body $j$ is:
    \begin{equation*}
        \mathbf{F}_{ij} = G \frac{m_i m_j}{|\mathbf{r}_j - \mathbf{r}_i|^3} (\mathbf{r}_j - \mathbf{r}_i)
    \end{equation*}
    Total force on body $i$ from all other bodies:
    \begin{equation*}
        \mathbf{F}_i = \sum_{\substack{j=1 \\ j \neq i}}^{N} \mathbf{F}_{ij}
    \end{equation*}
    Acceleration of body $i$:
    \begin{equation*}
        \mathbf{a}_i = \frac{\mathbf{F}_i}{m_i}
    \end{equation*}
\end{frame}

% Slide 8
\begin{frame}{The Softening Factor ($\epsilon$)}
    \textbf{The Problem of Singularities:}
    \begin{itemize}
        \item If two particles get extremely close, distance $|\mathbf{r}_{ij}| \to 0$.
        \item Force $\mathbf{F} \to \infty$.
        \item This breaks numerical integrators, causing particles to shoot off at infinite velocities.
    \end{itemize}
    \textbf{The Solution:} We introduce a small softening length $\epsilon$.
    \begin{equation*}
        \mathbf{F}_{ij} = G \frac{m_i m_j}{(|\mathbf{r}_{ij}|^2 + \epsilon^2)^{3/2}} \mathbf{r}_{ij}
    \end{equation*}
\end{frame}

% Slide 9
\begin{frame}{Initial Conditions}
    \textbf{The "Galaxy Disk" Model:}
    \begin{itemize}
        \item \textbf{Core:} One super-massive central particle ($M_{core} = N$).
        \item \textbf{Satellites:} $N-1$ particles initialized at random angles $\phi \in [0, 2\pi]$ and radii $R \in [1, 5]$.
        \item \textbf{Circular Velocity:} Initialized perpendicular to the radius.
        \begin{equation*}
            v = \sqrt{\frac{G \cdot M_{core}}{R}}
        \end{equation*}
        \item Ensures a stable, rotating disk configuration instead of immediate collapse.
    \end{itemize}
\end{frame}

% Slide 10
\begin{frame}{Unit Systems}
    \begin{itemize}
        \item \textbf{SI Units} ($kg, meters, seconds$) cause severe floating-point issues in astrophysics.
        \item Earth's mass is $\sim10^{24}$ kg, sizes are $10^{11}$ m.
        \item We use \textbf{Natural Dimensionless Units}:
        \begin{itemize}
            \item Gravitational Constant $G = 1.0$
            \item Satellite masses $\sim 1.0$
            \item Lengths mapped to Astronomical Units (AU)
        \end{itemize}
        \item Prevents numeric overflow and simplifies debugging.
    \end{itemize}
\end{frame}

\section{Direct Approach \& OpenMP}
% Slide 11
\begin{frame}{Algorithm: Direct $\mathcal{O}(N^2)$ Approach}
    \begin{itemize}
        \item Iterate through every body $i$.
        \item For every body $i$, iterate through every other body $j$.
        \item Calculate $\mathbf{F}_{ij}$.
        \item Highly computationally intensive, but perfectly accurate.
    \end{itemize}
\end{frame}

% Slide 12
\begin{frame}{Optimization: Newton's Third Law}
    \begin{itemize}
        \item Every action has an equal and opposite reaction: $\mathbf{F}_{ji} = -\mathbf{F}_{ij}$.
        \item We only need to compute unique pairs.
        \item Inner loop runs from $j = i + 1$ to $N$.
        \item Cuts total floating point operations exactly in half!
    \end{itemize}
\end{frame}

% Slide 13
\begin{frame}{Parallelizing with OpenMP}
    \begin{itemize}
        \item The outer loop over body $i$ can be distributed across multi-core CPUs.
        \item OpenMP creates a thread-pool (e.g., 16 hardware threads).
        \item Pragma directive: \texttt{\#pragma omp parallel for schedule(dynamic)}
        \item Why dynamic schedule? The inner loop gets smaller as $i$ grows, meaning early iterations take longer. Dynamic scheduling perfectly load-balances the threads.
    \end{itemize}
\end{frame}

% Slide 14
\begin{frame}[fragile]{The Data Race Problem}
    Because we use Newton's 3rd Law, Thread $T_1$ processing body $i$ updates body $j$'s force.
    At the \textit{exact same time}, Thread $T_2$ processing body $k$ might also try to update body $j$'s force.
    
    \begin{block}{Consequence}
        A memory race condition. Forces get overwritten, physics are corrupted.
    \end{block}
\end{frame}

% Slide 15
\begin{frame}{Solution: Thread-Local Reduction}
    \begin{itemize}
        \item We provide \textit{each parallel thread} its own fully initialized $\mathcal{O}(N)$ array of forces.
        \item Thread $T$ only writes to \texttt{localForce[T][j]}.
        \item Absolutely zero memory contention during the heavy $O(N^2)$ loop.
        \item Afterwards, a blazing-fast $\mathcal{O}(N \times Threads)$ sequential reduction merges all thread-local arrays into the global forces array.
    \end{itemize}
\end{frame}

\section{Barnes-Hut Algorithm}
% Slide 16
\begin{frame}{The Scaling Wall of $\mathcal{O}(N^2)$}
    \begin{itemize}
        \item Even with Newton's 3rd law and 16-core OpenMP parallelism, the Direct approach becomes physically un-runnable past $\sim 50,000$ particles.
        \item We need an algorithmic breakthrough, not just more hardware.
        \item Introduced in 1986 by Josh Barnes and Piet Hut.
    \end{itemize}
\end{frame}

% Slide 17
\begin{frame}{Introduction to Barnes-Hut}
    \begin{itemize}
        \item \textbf{Core Philosophy:} If a cluster of stars is sufficiently far away, we can approximate the entire cluster as a single massive point-particle located at the cluster's Center of Mass.
        \item Reduces complexity from checking $N$ bodies to traversing a tree in $\log N$ steps.
        \item Total complexity drops from $\mathcal{O}(N^2)$ to $\mathcal{O}(N \log N)$.
    \end{itemize}
\end{frame}

% Slide 18
\begin{frame}{The Octree Data Structure}
    \begin{itemize}
        \item A 3D spatial partitioning tree (the 3D equivalent of a Quadtree).
        \item The entire universe is a root cube.
        \item If a cube contains more than one particle, it is recursively subdivided into 8 smaller child octants.
        \item Leaves contain exactly 0 or 1 particle.
        \item Internal nodes store exactly the total mass and Center of Mass for all children below them.
    \end{itemize}
\end{frame}

% Slide 19
\begin{frame}{The Multipole Acceptance Criterion (MAC)}
    How do we decide if a node is "far enough" to be approximated?
    \begin{equation*}
        \frac{s}{d} < \theta
    \end{equation*}
    \begin{itemize}
        \item $s =$ width of the octree cell.
        \item $d =$ distance from the particle to the cell's center of mass.
        \item $\theta =$ threshold parameter (typically 0.5).
        \item If true, compute Force against the node. If false, recurse deeper into the children.
    \end{itemize}
\end{frame}

% Slide 20
\begin{frame}{Parallelizing Barnes-Hut}
    \begin{itemize}
        \item \textbf{Tree Construction:} Performed sequentially on 1 thread. It is only $\mathcal{O}(N \log N)$ and involves vast memory allocation, making thread-safety expensive.
        \item \textbf{Force Traversal:} Performed in parallel! 
        \item Since the tree is \textit{read-only} during force computation, each particle can traverse the tree simultaneously on different threads without any locks or race conditions.
    \end{itemize}
\end{frame}

\section{Time Integration}
% Slide 21
\begin{frame}{Time Integration: Why not Euler?}
    \begin{itemize}
        \item Euler integration: $\mathbf{x}_{new} = \mathbf{x}_{old} + \mathbf{v} \cdot dt$.
        \item It is a first-order method.
        \item It does not conserve energy! Gravity causes artificial outward spirals resulting in the galaxy exploding over time.
    \end{itemize}
\end{frame}

% Slide 22
\begin{frame}{The Leapfrog (Velocity-Verlet) Integrator}
    \textbf{Kick-Drift-Kick Scheme:}
    \begin{enumerate}
        \item \textbf{Kick:} Accelerate velocity by half a step ($dt/2$).
        \item \textbf{Drift:} Update position using the new half-step velocity.
        \item \textbf{Forces:} Re-calculate gravitational forces at new positions.
        \item \textbf{Kick:} Accelerate velocity by the second half step ($dt/2$).
    \end{enumerate}
    Calculates velocity at half-integer time intervals offset from positions.
\end{frame}

% Slide 23
\begin{frame}{Symplectic Properties}
    \begin{itemize}
        \item Leapfrog is a \textbf{Symplectic Integrator}.
        \item It preserves the phase-space volume of the system.
        \item \textbf{Real-world meaning:} Energy oscillates slightly step-to-step, but the long-term energy error is strictly bounded. Planets maintain stable, closed orbital ellipses forever instead of drifting.
    \end{itemize}
\end{frame}

\section{Results \& Visualizations}
% Slide 24
\begin{frame}{Software Architecture \& Toolchain}
    \begin{itemize}
        \item \textbf{Engine:} Modern C++17.
        \item \textbf{Build System:} CMake for cross-platform Multi-Target compilation.
        \item \textbf{Parallelism API:} OpenMP via \texttt{omp.h}.
        \item \textbf{Data Hand-off:} Computations write state-vectors efficiently to CSV files per-step.
        \item \textbf{Visualizer:} Python, Pandas, Matplotlib 3D, Headless Agg backend for GIF generation.
    \end{itemize}
\end{frame}

% Slide 25
\begin{frame}{Visualizing Trajectories}
    \begin{itemize}
        \item An automated pipeline reads \texttt{trajectories.csv}.
        \item Generates a 3-D animated GIF showing particle tracks over time.
        \item Generates a 2-D XY "Spaghetti-plot" of the disk.
        \item Shows clear galaxy stability, confirming appropriate Initial Conditions and Integrator capability.
    \end{itemize}
\end{frame}

% Slide 26
\begin{frame}{Performance Benchmark Results}
    Tests run across varying $N$ counts using 16 Threads on 500 Steps:
    \begin{table}[]
        \centering
        \begin{tabular}{|l|l|l|}
        \hline
        \textbf{Bodies (N)} & \textbf{Direct Time} & \textbf{Barnes-Hut Time} \\ \hline
        50                  & 1.23 s               & 0.97 s                   \\ \hline
        200                 & 1.81 s               & 1.66 s                   \\ \hline
        800                 & 5.36 s               & 5.00 s                   \\ \hline
        1600                & 9.75 s               & 9.68 s                   \\ \hline
        \end{tabular}
    \end{table}
\end{frame}

% Slide 27
\begin{frame}{Analyzing the Scaling Graph}
    \begin{itemize}
        \item Why is Direct comparable to Barnes-Hut at $N < 1600$?
        \item Array-based direct access is highly cache-friendly. Vector units in the CPU breeze through contiguous floating point arrays.
        \item Barnes-Hut requires pointer-chasing through tree nodes, which incurs significant CPU cache-misses and memory latency.
        \item But as $N \to 10000+$, mathematics takes over, and $O(N \log N)$ easily overcomes cache-miss latency.
    \end{itemize}
\end{frame}

\section{Conclusion}
% Slide 28
\begin{frame}{Conclusion}
    \begin{itemize}
        \item Successfully modeled complex gravitational dynamics.
        \item Showcased the extreme importance of Algorithm choice ($O(N^2)$ vs $O(N \log N)$) alongside Hardware choice (OpenMP Multi-threading).
        \item Proved that Leapfrog integration is necessary for energetic stability in Hamiltonian physics systems.
    \end{itemize}
\end{frame}

% Slide 29
\begin{frame}{Future Work}
    \begin{itemize}
        \item \textbf{GPGPU Acceleration:} Porting the nested $\mathcal{O}(N^2)$ loops to CUDA architectures to utilize 5000+ GPU cores rather than 16 CPU cores.
        \item \textbf{Collisions:} Implementing inelastic collision merging for celestial body formations.
        \item \textbf{Distributed Computing:} Implementing MPI (Message Passing Interface) to run the simulation across clusters of multiple computers.
    \end{itemize}
\end{frame}

% Slide 30
\begin{frame}
    \centering
    \Huge \textbf{Questions?} \\
    \vspace{1cm}
    \large Thank you for your time.
\end{frame}

\end{document}
